\begin{appendices}


\chapter{Joint distributions as transport plans}


The key idea is that a \textbf{joint distribution} does more than describe two random variables simultaneously---it specifies \textbf{how probability mass from one distribution is matched to probability mass in another}. That matching is exactly what is meant by a \textbf{transport plan} in optimal transport.\\

Suppose you have

\begin{itemize}
    \item a real data distribution $\mu$ on $\mathcal{X}$,
    \item a generated distribution $\nu$ on $\mathcal{Y}$.
\end{itemize}

By themselves, these distributions tell you only \textbf{where the mass is located}. They do \textbf{not} tell you which part of the mass in $\mu$ should be transported to which part of the mass in $\nu$.

A \textbf{joint distribution}

\[
\pi(x,y)
\]

defined on the product space

\[
\mathcal{X}\times\mathcal{Y}
\]

provides exactly this missing information.\\


A transport plan specifies the amount of mass moved from every point $x$ to every point $y$.

More precisely,

\[
\pi(A\times B)
\]

represents the amount of probability mass originally in the subset $A\subseteq\mathcal{X}$ that is transported to the subset $B\subseteq\mathcal{Y}$.

Equivalently, if $\pi$ has a density, then

\[
\pi(x,y)\,dx\,dy
\]

is the infinitesimal amount of mass transported from a neighborhood around $x$ to a neighborhood around $y$.

Thus every pair $(x,y)$ corresponds to one possible transportation route, and the joint distribution tells us how much mass uses each route.\\

A transport plan cannot create or destroy mass.

Therefore,

\[
\int_{\mathcal Y}\pi(x,y)\,dy=\mu(x)
\]

states that all mass leaving $x$ equals the amount of mass originally located at $x$.

Similarly,

\[
\int_{\mathcal X}\pi(x,y)\,dx=\nu(y)
\]

states that all mass arriving at $y$ equals the amount of mass required at $y$.

These are exactly the marginal constraints

\[
\pi\in\Pi(\mu,\nu).
\]

They enforce \textbf{mass conservation}.

\section*{A simple discrete example}

Suppose

\[
\mu=(0.6,\,0.4)
\]

and

\[
\nu=(0.5,\,0.5).
\]

One admissible joint distribution is

\[
\pi=
\begin{pmatrix}
0.5 & 0.1\\
0 & 0.4
\end{pmatrix}.
\]

Its row sums are

\[
(0.6,\;0.4)=\mu,
\]

and its column sums are

\[
(0.5,\;0.5)=\nu.
\]

This matrix has a natural transportation interpretation:

\begin{itemize}
    \item transport $0.5$ units of mass from $x_1$ to $y_1$,
    \item transport $0.1$ units from $x_1$ to $y_2$,
    \item transport $0.4$ units from $x_2$ to $y_2$.
\end{itemize}

Nothing else is transported.

Notice that the matrix is simultaneously

\begin{itemize}
    \item a joint probability distribution,
    \item a feasible transport plan.
\end{itemize}

The two viewpoints are mathematically identical.

\section*{Connection to the Kantorovich formulation}

The Kantorovich problem minimizes the total transportation cost

\[
\inf_{\pi\in\Pi(\mu,\nu)}
\int_{\mathcal X\times\mathcal Y}
c(x,y)\,d\pi(x,y).
\]

Here,

\begin{itemize}
    \item $c(x,y)$ is the cost of transporting one unit of mass from $x$ to $y$,
    \item $d\pi(x,y)$ is the amount of mass transported along that route.
\end{itemize}

Therefore,

\[
c(x,y)\,d\pi(x,y)
\]

is the infinitesimal transportation cost, and integrating over all pairs $(x,y)$ gives the total work required to transform $\mu$ into $\nu$.

\section*{Connection to WGANs}

In WGANs, the real distribution $P_r$ and the generated distribution $P_g$ are generally known only through samples. The transport plan $\pi$ is \textbf{not computed explicitly}. Instead, the Kantorovich duality shows that the same optimal transport cost can be obtained by optimizing over 1-Lipschitz functions (the critic), avoiding the need to solve directly for $\pi$.

Nevertheless, conceptually, the transport plan remains the object that specifies \textbf{how probability mass from the real distribution is matched to probability mass in the generated distribution}. The critic can be viewed as evaluating the cost of this optimal matching without ever constructing the joint distribution itself.







\end{appendices}